\section{String Matching}
\label{sec:cses-string-matching}

\problemstatement{Given a text $t$ and a pattern $p$, count how many times
$p$ occurs in $t$ (occurrences may overlap).}

\begin{patternbox}
\emph{``Count exact occurrences of a pattern''} is solved in linear time
by the \textbf{Z-function} (or KMP). Concatenate $p+\#+t$ with a separator
$\#$ not in either string; a position in the $t$-part whose Z-value equals
$|p|$ is an occurrence.
\end{patternbox}

\subsection*{The Z-function}

The Z-function of a string $w$ gives, for each index $i$, the length of the
longest substring starting at $i$ that matches a prefix of $w$. Build
$w=p+\#+t$. Wherever $Z[i]=|p|$ and $i$ lies in the $t$ portion, the text
matches the pattern's full length at that spot -- an occurrence. The
separator $\#$ guarantees a match cannot span the boundary. The Z-array
itself is computed in $O(|w|)$ using a maintained ``Z-box'' -- the
rightmost match interval seen so far.

\begin{ideabox}[Reuse Previous Matches via the Z-box]
When computing $Z[i]$, if $i$ falls inside a previously found match
interval $[l,r]$, the value $Z[i-l]$ is already known and can seed
$Z[i]$, so each character is compared a constant number of times
amortised. This is what makes Z (and KMP) linear rather than quadratic.
\end{ideabox}

\subsection*{Algorithm}

\begin{enumerate}
  \item Form $w=p+\#+t$ and compute its Z-array.
  \item Count indices $i$ (in the $t$ region) with $Z[i]=|p|$.
  \item Output the count.
\end{enumerate}

\subsection*{Dry run}

\dryrun{$t=$``saippua'', $p=$``pp''; $w=$``pp\#saippua''.}

\begin{itemize}
  \item $|p|=2$. Scanning $w$, the substring ``pp'' inside $t$ starts at
        one position; there $Z=2$.
  \item Exactly one index has $Z=2$, so the pattern occurs once.
\end{itemize}

\subsection*{C++ implementation}

\begin{lstlisting}
#include <bits/stdc++.h>
using namespace std;

vector<int> zFunction(const string& w) {
    int n = w.size();
    vector<int> z(n, 0);
    z[0] = n;
    int l = 0, r = 0;
    for (int i = 1; i < n; i++) {
        if (i < r) z[i] = min(r - i, z[i - l]);
        while (i + z[i] < n && w[z[i]] == w[i + z[i]]) z[i]++;
        if (i + z[i] > r) { l = i; r = i + z[i]; }
    }
    return z;
}

int main() {
    string t, p;
    cin >> t >> p;
    string w = p + "#" + t;
    vector<int> z = zFunction(w);

    int count = 0, m = p.size();
    for (int i = m + 1; i < (int)w.size(); i++)   // t region
        if (z[i] == m) count++;
    cout << count << "\n";
    return 0;
}
\end{lstlisting}

\begin{complexitybox}
\textbf{Time Complexity.} $O(|p|+|t|)$ -- one linear Z-array pass.
\textbf{Space Complexity.} $O(|p|+|t|)$ for the combined string and its
Z-array.
\end{complexitybox}

\begin{takeawaybox}
Exact pattern counting reduces to ``where does a prefix of $p+\#+t$
reappear at full pattern length,'' which the Z-function answers in linear
time. KMP's prefix function gives an equivalent solution; keep both in
your toolkit as interchangeable linear matchers.
\end{takeawaybox}
